\scanpage{32}

\begin{samepage}
\tikz[remember picture,baseline=(lgstart.base)]{\node[inner sep=0pt](lgstart){1.};} Généralisation des th. de finitude formelle, et des th. de comparaison. Application : un th. d'existence

2. Le Théorème de finitude locale et ses applications.

\deleted{3. Fermés de Lefschetz-Grauert\unsure{:|.} Remarques\unsure{.|,}}

\deleted{4. \unsure{Premières|Propriétés} propriétés \unsure{des} \uline{Exemples} \unsure{fondamentaux|fondament}} \inserted{1}

3. Cas particulier 1 : \deleted{\unsure{Etude|Étude}} \inserted{théorie} \unsure{des|d'en} \deleted{\unsure{spa}} schémas locaux.

4. Cas particulier 2 : \deleted{\unsure{Etude|Étude}} \inserted{théorie} des schémas \inserted{algébriques} projectifs.

5. Fermés de Lefschetz-Grauert : Remarques, exemples.

% slanted
\margin{\tikz[remember picture,baseline=(lgnote.base)]{\node[inner sep=1pt,text width=1.9cm,draw,align=left](lgnote){utilise IV, \S\ sur les morphismes étales....};}}
6. Propriétés liées aux $\pi_0$, \tikz[remember picture,baseline=(lgpi.base)]{\node[draw,ellipse,inner sep=1pt](lgpi){$\pi_1$};}, Pic.

[7. Propriétés de permanence]

\tikz[remember picture,baseline=(lgeight.base)]{\node[draw,ellipse,inner sep=1pt](lgeight){8.};} Théorèmes de pureté, et théorèmes de Lefschetz pour le $\pi_1$.

\deleted{\textcircled{\unsure{9|1}}} \tikz[remember picture,baseline=(lgnine.base)]{\node[inner sep=0pt](lgnine){9};} \uline{Anneaux} \deleted{\unsure{locaux}} \uline{quasi-factoriels} et \uline{th. de Lefschetz pour Pic}.

\tikz[remember picture,overlay]{\draw[->] (lgnote.east) to[bend left=12] (lgpi.north west); \draw[->] (lgnote.south east) to[bend right=12] (lgeight.west); \draw ([xshift=-4pt]lgstart.north west) -- ([xshift=-12pt]lgstart.north west) -- ([xshift=-12pt]lgnine.south west) -- ([xshift=-4pt]lgnine.south west); \draw ([xshift=-12pt,yshift=7pt]lgnine.north west) -- ([xshift=-5pt,yshift=7pt]lgnine.north west);}
\end{samepage}

% archive sheet number: 31
